%%=============================================================================
%% Voorwoord
%%=============================================================================

\chapter*{\IfLanguageName{dutch}{Woord vooraf}{Preface}}%
\label{ch:voorwoord}

%% TODO:
%% Het voorwoord is het enige deel van de bachelorproef waar je vanuit je
%% eigen standpunt (``ik-vorm'') mag schrijven. Je kan hier bv. motiveren
%% waarom jij het onderwerp wil bespreken.
%% Vergeet ook niet te bedanken wie je geholpen/gesteund/... heeft

%%=============================================================================
%% Voorwoord
%%=============================================================================

Het onderwerp van deze bachelorproef, gericht op de ontwikkeling en evaluatie van een privacybewust spelling- en grammaticacontrole-toetsenbord voor Android, werd gedaan vanwege mijn interesse in de exponentiële evolutie van artificiële intelligentie in de afgelopen jaren en de bijbehorende beveiligingsaspecten. De mogelijkheid om AI toe te passen op een manier die zowel functioneel als veilig is, fascineerde mij en diende als een belangrijke drijfveer voor dit onderzoek.

\vspace{0.3cm}

Ik wil graag mijn oprechte dank uitspreken aan mijn promotor en co-promotor, de heer Jan Willem, voor zijn begeleiding en ondersteuning gedurende dit project. Zijn snelle reacties op mijn vragen en constante bijstand hebben me geholpen om tot dit eindresultaat te komen. Zijn waardevolle adviezen hebben een belangrijke bijdrage geleverd aan mijn onderzoek.

\vspace{0.3cm}

Daarnaast wil ik mijn waardering uitspreken voor de individuen die, ondanks de fysieke afstand, een grote steun voor mij zijn geweest. Hun aanmoediging om door te gaan heeft een cruciale rol gespeeld tijdens momenten van uitdaging. Speciale dank gaat uit naar Pieter Delobelle, de maker van RobBERT, voor het beantwoorden van mijn vragen over RobBERT, en Alex, de ontwikkelaar van Futo Keyboard, voor zijn inzichten in de ontwikkeling van een Android-toetsenbord.

\vspace{0.3cm}

De reis van deze bachelorproef was niet zonder uitdagingen. Er waren momenten waarin de complexiteit van het onderwerp overweldigend leek, maar door mezelf te verdiepen in de noodzakelijke AI-onderwerpen, wist ik deze obstakels te overwinnen. Dit proces heeft niet alleen mijn begrip van artificiële intelligentie verdiept, maar ook mijn passie voor het veld aangewakkerd.

\vspace{0.3cm}

Ik hoop dat de kennis die ik door dit project heb opgedaan, zal dienen als basis voor het creëren van innovatieve software in de toekomst. Mijn streven is om technologieën te ontwikkelen die miljoenen mensen wereldwijd kunnen helpen.

\vspace{0.3cm}

Met dank aan allen die op enige wijze hebben bijgedragen, presenteer ik deze bachelorproef met trots en hoop ik dat het een waardevolle bijdrage levert aan het domein van privacybewuste technologieën.

\vspace{1cm}

\begin{flushright}
    Roy Barneveld\\
    Gent\\
    20 Augustus 2024
\end{flushright}